\documentclass[../main.tex]{subfiles}
\usepackage{enumitem}
\newlist{romanlist}{enumerate}{1}
\setlist[romanlist]{label=(\roman*), leftmargin=2.2em, itemsep=2pt}
\begin{document}

Having described the proposed three-tier hybrid system in
Chapter~3, this chapter evaluates it empirically. The aim
is twofold: to quantify how well each individual tier performs, and to show
that the smart hybrid fusion improves upon every single tier on the metric
that matters most for intrusion detection. Section~\ref{sec:eval-params}
defines the evaluation parameters, that is, the metrics used and the
composition of the evaluation set. Section~\ref{sec:sim-method} details the
simulation method: the dataset preparation, the baselines selected for
comparison and the rationale behind them, the algorithm hyperparameters, and
the experimental procedure. The results are then reported in three parts:
the performance of the individual detection tiers
(Section~\ref{sec:res-individual}), the performance of the hybrid fusion
strategies (Section~\ref{sec:res-hybrid}), and an error analysis based on
confusion matrices (Section~\ref{sec:res-confusion}). Each result is
presented with the corresponding tables and figures, accompanied by a
detailed discussion that compares the methods and explains the observed
behaviour.

% =====================================================================
\section{Evaluation Parameters}
\label{sec:eval-params}

The system performs binary classification, in which a request flagged as an
attack is treated as the \emph{positive} class and a benign request as the
\emph{negative} class. Its performance is therefore measured through the
confusion matrix and the metrics derived from it. The confusion matrix, whose
layout is shown in Table~\ref{tab:cm}, cross-tabulates the predicted class
against the true class and yields four counts. The number of true positives
($TP$) is the number of attacks that are correctly detected; the number of
true negatives ($TN$) is the number of benign requests that are correctly
passed; the number of false positives ($FP$) is the number of benign requests
that are wrongly flagged as attacks (the \emph{false alarms}); and the number
of false negatives ($FN$) is the number of attacks that are missed (the
\emph{missed detections}). The two diagonal cells of the matrix therefore
count the correct decisions, while the two off-diagonal cells count the two
kinds of error.

\begin{table}[ht]
  \centering
  \caption{Layout of the binary confusion matrix used in this chapter.}
  \label{tab:cm}
  \renewcommand{\arraystretch}{1.35}
  \setlength{\tabcolsep}{10pt}
  \footnotesize
  \begin{tabular}{|l|c|c|}
    \hline
    & \textbf{Predicted: Attack} & \textbf{Predicted: Clean} \\
    \hline
    \textbf{Actual: Attack} & $TP$ (correct) & $FN$ (missed attack) \\
    \hline
    \textbf{Actual: Clean} & $FP$ (false alarm) & $TN$ (correct) \\
    \hline
  \end{tabular}
\end{table}

From these four counts, four metrics are reported for every configuration.
\emph{Precision} is the fraction of the requests flagged as attacks that are
genuinely attacks; it is high when there are few false alarms, and is defined
as
\begin{equation}
  \text{Precision} = \frac{TP}{TP + FP}.
  \label{eq:precision}
\end{equation}
\emph{Recall}, also called the detection rate or the true-positive rate, is
the fraction of the actual attacks that the system manages to detect; it is
high when there are few missed detections, and is defined as
\begin{equation}
  \text{Recall} = \frac{TP}{TP + FN}.
  \label{eq:recall}
\end{equation}
For reference, \emph{accuracy} is the fraction of all decisions that are
correct,
\begin{equation}
  \text{Accuracy} = \frac{TP + TN}{TP + TN + FP + FN}.
  \label{eq:accuracy}
\end{equation}
Precision and recall capture complementary kinds of error and are usually in
tension: a detector can trivially achieve perfect recall by flagging every
request (at the cost of precision), or near-perfect precision by flagging only
the most obvious attacks (at the cost of recall). To summarise both in a
single number, the F-measure takes their weighted harmonic mean. The general
$F_\beta$ score, in which the parameter $\beta$ sets how many times more
important recall is than precision, is
\begin{equation}
  F_\beta = (1 + \beta^2)\,\frac{\text{Precision} \cdot \text{Recall}}
                 {\beta^2 \cdot \text{Precision} + \text{Recall}}.
  \label{eq:fbeta}
\end{equation}
Setting $\beta = 1$ weights precision and recall equally and gives the
familiar F1 score, the plain harmonic mean,
\begin{equation}
  F_1 = \frac{2 \cdot \text{Precision} \cdot \text{Recall}}
             {\text{Precision} + \text{Recall}}.
  \label{eq:f1}
\end{equation}
Setting $\beta = 2$ weights recall twice as heavily as precision and gives the
F2 score, which after substituting $\beta = 2$ into
Equation~\ref{eq:fbeta} simplifies to
\begin{equation}
  F_2 = \frac{5 \cdot \text{Precision} \cdot \text{Recall}}
             {4 \cdot \text{Precision} + \text{Recall}}.
  \label{eq:f2}
\end{equation}

In this thesis the F2 score of Equation~\ref{eq:f2} is treated as the primary
metric of merit. The reason is rooted in the cost asymmetry of the intrusion
detection task: a missed attack (a false negative) may lead to a data breach
with severe consequences, whereas a false alarm (a false positive) merely
costs an analyst a few minutes of investigation. Because the two errors are
not equally damaging, a metric that weights recall above precision reflects
the operational priority better than the balanced F1 score does. Accuracy,
defined in Equation~\ref{eq:accuracy}, is reported only for completeness and
is not used as a headline metric, because on a near-balanced dataset it is
less discriminative than the precision--recall family and can mask a poor
recall.

All headline numbers in this chapter are computed on a single, fixed
evaluation set whose composition is given in Table~\ref{tab:evalset}. This
set is the held-out $30\%$ split of the balanced corpus; after log parsing,
$28{,}767$ lines were successfully evaluated, of which $13{,}643$ are
attacks and $15{,}124$ are clean, giving an attack ratio of $47.43\%$.
Because the two classes are close to balanced, both precision and recall are
meaningful and neither is trivially inflated by class skew.

\begin{table}[ht]
  \centering
  \caption{Composition of the primary evaluation set.}
  \label{tab:evalset}
  \renewcommand{\arraystretch}{1.35}
  \setlength{\tabcolsep}{10pt}
  \footnotesize
  \begin{tabular}{|l|r|}
    \hline
    \textbf{Quantity} & \textbf{Value} \\
    \hline
    Evaluated log lines & 28{,}767 \\
    \hline
    Attack lines & 13{,}643 \\
    \hline
    Clean lines & 15{,}124 \\
    \hline
    Attack ratio & 47.43\% \\
    \hline
  \end{tabular}
\end{table}

% =====================================================================
\section{Simulation Method}
\label{sec:sim-method}

\subsection{Dataset preparation}
\label{sec:sim-data}

The corpus combines three sources of HTTP traffic, as introduced in
Section~\ref{sec:scope}: the CSIC~2010 dataset, real attack payloads drawn
from the public PayloadsAllTheThings repository, and automatically generated
synthetic traffic. All three sources are assembled by a single, fully
reproducible build script,
\texttt{build\_\allowbreak final\_\allowbreak dataset\_\allowbreak comprehensive.py}.
Before balancing, the pool contains $111{,}065$ log lines whose provenance is
detailed in Table~\ref{tab:sources}: the two CSIC subsets supply realistic
e-commerce traffic, while the synthetic generator adds modern attack classes
and a controlled variety of benign endpoints.

\begin{table}[ht]
  \centering
  \caption{Composition of the raw corpus before balancing (the
    $111{,}065$-line pool assembled by the build script).}
  \label{tab:sources}
  \renewcommand{\arraystretch}{1.25}
  \setlength{\tabcolsep}{8pt}
  \footnotesize
  \begin{tabular}{|l|l|r|r|}
    \hline
    \textbf{Source} & \textbf{Class} & \textbf{Lines} & \textbf{Share} \\
    \hline
    CSIC~2010 (clean)   & Clean  & $36{,}000$  & $32.4\%$ \\
    CSIC~2010 (attack)  & Attack & $25{,}065$  & $22.6\%$ \\
    Synthetic (clean)   & Clean  & $25{,}000$  & $22.5\%$ \\
    Synthetic (attack)  & Attack & $25{,}000$  & $22.5\%$ \\
    \hline
    \textbf{Total}      &        & $111{,}065$ & $100\%$ \\
    \hline
  \end{tabular}
\end{table}

The synthetic \emph{attack} traffic is produced as follows. The generator
draws on a curated set of $448$ real attack payloads organised into six
categories --- SQL injection, cross-site scripting, XML external entity
injection, remote command execution, CRLF injection, and LDAP injection ---
taken from the PayloadsAllTheThings repository. For each of the $25{,}000$
synthetic attack lines it picks a category and a payload at random and then
applies a random obfuscation chosen from the raw payload, a single
URL-encoding, a double URL-encoding, and a space-to-\texttt{\%20}
substitution; this deliberately exercises the multi-layer decoder of
Chapter~3 and prevents the models from memorising a fixed surface form. The
payload is embedded either in the query string of a \texttt{GET} request
(sixty per cent of the lines) or in the body of a \texttt{POST} request
(forty per cent), and the line is completed with a random client address, a
timestamp spread uniformly over a thirty-day window, a status code drawn from
$\{200,403,500\}$, a response size, and a realistic User-Agent string. A
hidden marker is appended to the User-Agent so that the ground-truth label of
every synthetic line is exact rather than inferred.

The synthetic \emph{clean} traffic is generated by the same script in a
parallel manner. Each of the $25{,}000$ benign lines targets a random path
from a catalogue of typical application endpoints (for example
\texttt{/index.html}, \texttt{/api/users}, \texttt{/login}, \texttt{/search},
static assets, and \texttt{/dashboard}); it is issued as a \texttt{GET} with
ordinary pagination parameters (eighty per cent) or as a \texttt{POST} with
ordinary form parameters (twenty per cent), again with a randomised address,
timestamp, size, and User-Agent, and always returning status~$200$. Because
these lines are structurally indistinguishable from genuine benign traffic
but carry no marker, they teach the models the \emph{shape} of normal requests
rather than any artefact peculiar to a single dataset.

Balancing the corpus to exactly $50\%$ attack and $50\%$ clean is a
deliberate choice. On a heavily imbalanced set a trivial classifier can reach
a high accuracy merely by always predicting the majority class, which would
mask its true discriminative ability; an even split removes this class-prior
shortcut and makes precision and recall reflect genuine detection skill. The
attack pool ($50{,}065$ lines) is the limiting class, so the clean pool
($61{,}000$ lines) is randomly down-sampled to the same size, giving a
balanced corpus of $100{,}130$ lines. This corpus is shuffled and split
$70/30$ into $70{,}091$ training and $30{,}039$ evaluation lines. Both splits
are then passed through the parsing pipeline of Chapter~3 --- each line is
parsed, its URL decoded over up to three layers, and converted into the shared
$22$-feature vector --- and the small number of lines that cannot be parsed
are discarded. This leaves $66{,}984$ usable labelled lines for the supervised
tier (Tier~2), of which $32{,}043$ are attacks and $34{,}941$ are clean, and an
evaluation set of $28{,}767$ lines (Table~\ref{tab:evalset}). Following the
one-class principle, the unsupervised tier (Tier~3) is trained on the
$34{,}941$ clean training lines only. Because every label is derived from the
embedded marker rather than estimated, the evaluation labels are exact;
static-resource requests are whitelisted out of the learning tiers, as in
production, so that the reported numbers reflect the system's real operating
behaviour.

\subsection{Baselines and rationale}
\label{sec:sim-baselines}

To place the proposed hybrid in context, it is compared against a set of
baselines that span the design space. The baselines, and the reason each was
chosen, are the following:

\begin{romanlist}
  \item \textbf{Tier~1 regex rules}, as a pure misuse-detection baseline.
        It represents the classical signature approach and establishes the
        precision ceiling and the recall floor of the system.
  \item \textbf{Logistic Regression}, as a linear supervised baseline against
        the non-linear Random Forest. It was chosen because it is the model
        reported as best by Pham et al.~\cite{pham2016} on CSIC 2010, which
        makes the comparison directly relevant to the related work.
  \item \textbf{Isolation Forest} and \textbf{One-Class SVM}, as alternative
        one-class anomaly detectors against the Local Outlier Factor. They
        were chosen so that the choice of the operational Tier~3 model rests
        on an explicit comparison rather than an assumption.
  \item \textbf{Simple Voting} (alert if any tier fires, that is, the logical
        OR of the three tiers), as a naive fusion baseline against the
        proposed smart consensus. It isolates the contribution of the
        consensus logic from the mere act of combining tiers.
\end{romanlist}

This selection lets the evaluation answer three distinct questions: which
supervised model is best, which unsupervised model is best, and whether the
smart consensus genuinely improves on both the single tiers and a naive
fusion.

\subsection{Algorithm hyperparameters}
\label{sec:sim-hyper}

The hyperparameters of every algorithm are fixed in advance and summarised in
Table~\ref{tab:hyper}. The two supervised models use balanced class weights
to counteract any residual class imbalance during fitting. The Random Forest
and the Isolation Forest use $200$ estimators. The Local Outlier Factor uses
twenty neighbours in novelty mode so that it can score unseen requests. The
One-Class SVM uses a radial-basis-function kernel. The models that are
sensitive to feature scale (Logistic Regression and One-Class SVM) are
preceded by a standardisation step, whereas the tree-based models are not.
Every estimator is given a fixed random seed, which makes training and
inference deterministic.

\begin{table}[ht]
  \centering
  \caption{Algorithm hyperparameters used in the experiments.}
  \label{tab:hyper}
  \renewcommand{\arraystretch}{1.35}
  \setlength{\tabcolsep}{8pt}
  \footnotesize
  \begin{tabular}{|l|p{8.2cm}|}
    \hline
    \textbf{Model} & \textbf{Hyperparameters} \\
    \hline
    Random Forest (Tier 2) & 200 trees, unlimited depth, balanced class
    weights, fixed random seed \\
    \hline
    Logistic Regression (Tier 2) & lbfgs solver, up to 2000 iterations,
    balanced class weights, with feature standardisation \\
    \hline
    Local Outlier Factor (Tier 3) & 20 neighbours, novelty mode,
    contamination $=0.50$, with feature standardisation \\
    \hline
    Isolation Forest (Tier 3) & 200 trees, contamination $=0.50$, fixed
    random seed, with feature standardisation \\
    \hline
    One-Class SVM (Tier 3) & RBF kernel, $\nu = 0.60$, scale gamma, with
    feature standardisation \\
    \hline
  \end{tabular}
\end{table}

\subsection{Experimental procedure}
\label{sec:sim-procedure}

The experiment evaluates each configuration on the full evaluation set by a
single exact per-sample pass: for every one of the $28{,}767$ lines, the
configuration's decision is computed and tallied into the confusion matrix,
from which the four metrics are derived. Because all estimators are seeded
and the evaluation visits every sample exactly once (rather than sampling a
subset), the procedure is deterministic; repeating it yields identical
numbers, so no averaging over multiple runs is required and no error bars are
reported. In total, six single-model configurations (regex, Random Forest,
Logistic Regression, Isolation Forest, One-Class SVM, Local Outlier Factor)
and the hybrid fusion strategies (simple voting and the smart consensus) are
evaluated under identical conditions, which guarantees that the comparison is
fair and that any difference is attributable to the method rather than to the
data or the protocol.

% =====================================================================
\section{Performance of the Individual Detection Tiers}
\label{sec:res-individual}

The first experiment measures each detector in isolation. Table~\ref{tab:individual}
reports the four metrics for the regex baseline (Tier~1), the two supervised
models (Tier~2), and the three one-class models (Tier~3).
Figure~\ref{fig:metrics} visualises the same data as four grouped bar charts,
and Figure~\ref{fig:f1rank} ranks the models by F1 score.

\begin{table}[ht]
  \centering
  \caption{Performance of the individual detectors on the evaluation set
           (all values in \%).}
  \label{tab:individual}
  \renewcommand{\arraystretch}{1.35}
  \setlength{\tabcolsep}{9pt}
  \footnotesize
  \begin{tabular}{|l|l|r|r|r|r|}
    \hline
    \textbf{Detector} & \textbf{Type} & \textbf{Prec.} & \textbf{Recall} &
    \textbf{F1} & \textbf{F2} \\
    \hline
    Regex rules (Tier 1) & Signature & 99.61 & 24.07 & 38.77 & 28.37 \\
    \hline
    Random Forest (Tier 2) & Supervised & 94.80 & 94.76 & \textbf{94.78} & 94.77 \\
    \hline
    Logistic Regression (Tier 2) & Supervised & 85.82 & 74.90 & 79.99 & 76.86 \\
    \hline
    Isolation Forest (Tier 3) & Unsupervised & 63.02 & 93.14 & 75.17 & 85.01 \\
    \hline
    One-Class SVM (Tier 3) & Unsupervised & 62.27 & 91.38 & 74.07 & 83.57 \\
    \hline
    Local Outlier Factor (Tier 3) & Unsupervised & 88.28 & 95.13 & \textbf{91.57} & 93.67 \\
    \hline
  \end{tabular}
\end{table}

\begin{figure}[ht]
  \centering
  \includegraphics[width=\textwidth]{Figure/comprehensive_metrics.png}
  \caption{Precision, recall, F1, and F2 of every model, shown as four
           grouped bar charts. Green bars are supervised models, blue bars
           are unsupervised models, and red bars are the hybrid
           configurations (discussed in Section~\ref{sec:res-hybrid}).}
  \label{fig:metrics}
\end{figure}

\begin{figure}[ht]
  \centering
  \includegraphics[width=0.85\textwidth]{Figure/f1_ranking.png}
  \caption{Ranking of the detectors by F1 score. The Random Forest leads,
           closely followed by the smart consensus and the Local Outlier
           Factor.}
  \label{fig:f1rank}
\end{figure}

Several observations follow from Table~\ref{tab:individual} and
Figures~\ref{fig:metrics} and~\ref{fig:f1rank}. First, the regex baseline
behaves exactly as the misuse-detection paradigm predicts: its precision is
the highest of all detectors at $99.61\%$, meaning that almost everything it
flags is truly an attack, but its recall is only $24.07\%$, meaning it misses
roughly three out of four attacks. Consequently its F1 ($38.77\%$) and F2
($28.37\%$) are the lowest in the table. This confirms that a signature-only
system, while precise, is inadequate as a standalone detector, and motivates
the learning tiers.

Second, among the supervised models the Random Forest is decisively better
than Logistic Regression, achieving an F1 of $94.78\%$ against $79.99\%$. The
gap is largest in recall ($94.76\%$ versus $74.90\%$), which indicates that
the linear decision boundary of Logistic Regression cannot capture the
non-linear interactions among the $22$ features as well as the ensemble of
trees. This result is also consistent with, yet distinct from, Pham et
al.~\cite{pham2016}, who found Logistic Regression best on a high-dimensional
TF--IDF representation; on the compact $22$-feature representation used here,
the non-linear Random Forest is clearly superior, which justifies its choice
as the primary Tier~2 model.

Third, among the one-class models the Local Outlier Factor is the clear
winner, with an F1 of $91.57\%$ compared with $75.17\%$ for the Isolation
Forest and $74.07\%$ for the One-Class SVM. Interestingly, all three reach a
high recall (above $91\%$), but the Isolation Forest and the One-Class SVM
pay for it with a low precision (around $62$--$63\%$), whereas the Local
Outlier Factor keeps precision high at $88.28\%$. The reason is that the
density-based, locally adaptive criterion of the Local Outlier Factor draws a
tighter boundary around the normal manifold than the global partitioning of
the Isolation Forest or the single margin of the One-Class SVM, so it admits
far fewer clean requests as false anomalies. This is precisely why the Local
Outlier Factor was selected as the operational Tier~3 model.

\subsection{Feature importance analysis of the Random Forest}
\label{sec:res-featimp}

Because the Random Forest is the strongest single detector, it is worth asking
\emph{which} of the $22$ features drive its decisions. A tree ensemble exposes
this directly through the Gini importance of each feature, that is, the total
reduction in node impurity attributable to splits on that feature, averaged
over all the trees and normalised to sum to one. Figure~\ref{fig:featimp}
ranks the $22$ features by this score and Table~\ref{tab:featimp} lists the
leading contributors together with their cumulative share. This analysis
serves two purposes: it answers the practical question of which signals matter
most, and it provides empirical evidence that the hand-crafted feature set of
Chapter~3 was designed for good reason.

\begin{figure}[ht]
  \centering
  \includegraphics[width=0.92\textwidth]{Figure/feature_importance.png}
  \caption{Gini importance of the $22$ features in the Random Forest (Tier~2),
           ranked from highest to lowest. Bars are coloured by magnitude:
           dominant ($\geq 10\%$), supporting ($1$--$10\%$), and marginal
           ($< 1\%$).}
  \label{fig:featimp}
\end{figure}

\begin{table}[ht]
  \centering
  \caption{Leading Random Forest features by Gini importance, with cumulative
           share. The remaining fourteen features together account for the
           final $15.2\%$.}
  \label{tab:featimp}
  \renewcommand{\arraystretch}{1.25}
  \setlength{\tabcolsep}{8pt}
  \footnotesize
  \begin{tabular}{|c|l|r|r|}
    \hline
    \textbf{Rank} & \textbf{Feature} & \textbf{Importance} & \textbf{Cumulative} \\
    \hline
    1 & \texttt{entropy}                 & $17.9\%$ & $17.9\%$ \\
    2 & \texttt{longest\_param\_value}    & $16.9\%$ & $34.9\%$ \\
    3 & \texttt{len\_url}                 & $14.4\%$ & $49.2\%$ \\
    4 & \texttt{param\_value\_max\_entropy} & $12.5\%$ & $61.8\%$ \\
    5 & \texttt{path\_depth}              & $7.2\%$  & $69.0\%$ \\
    6 & \texttt{risk\_ratio}              & $6.3\%$  & $75.2\%$ \\
    7 & \texttt{status\_code\_indicator}  & $4.8\%$  & $80.0\%$ \\
    8 & \texttt{raw\_risk\_count}         & $4.8\%$  & $84.8\%$ \\
    \hline
    \multicolumn{2}{|l|}{Remaining $14$ features} & $15.2\%$ & $100\%$ \\
    \hline
  \end{tabular}
\end{table}

The ranking is highly interpretable. The four dominant features, which
together account for $61.8\%$ of the total importance, all measure
\emph{randomness} or \emph{length}: the Shannon entropy of the request
(\texttt{entropy}), the length of the longest parameter value
(\texttt{longest\_param\_value}), the overall URL length (\texttt{len\_url}),
and the maximum entropy among the parameter values
(\texttt{param\_value\_max\_entropy}). This matches the intrinsic nature of
injection attacks: a SQL-injection, cross-site-scripting, or XXE payload is
typically far longer and far more irregular than a normal parameter, so these
four signals separate the two classes most sharply. The next tier of features
--- the depth of the path (\texttt{path\_depth}), the density of risky tokens
(\texttt{risk\_ratio}), the response status (\texttt{status\_code\_indicator}),
the raw count of risky tokens (\texttt{raw\_risk\_count}), and the proportion
of encoded characters (\texttt{encoding\_ratio}) --- adds structural, lexical
and behavioural corroboration. Together, just \emph{eight} of the $22$ features
already explain $84.8\%$ of the model's decisions, which shows that the
representation is compact and that its features are complementary rather than
redundant.

It is equally important to interpret the features with \emph{low} importance,
because a low Gini score does not by itself mean a feature is useless. Three
features score essentially zero on this corpus, and each illustrates a
different reason. The first, \texttt{unknown\_param\_ratio}, measures how many
of a request's parameter names are absent from a per-endpoint vocabulary of
names previously seen for that exact path; it is the signal behind the Tier~1
parameter-tampering rule. In the training configuration used here that
vocabulary was not populated, so the feature was constant and the Random
Forest could not split on it. Far from being superfluous, it is expected to
become discriminative once the system is deployed in front of a specific
application and the vocabulary is built from that application's own clean
logs, at which point it can flag tampering attacks that reuse legitimate-looking
values under unexpected parameter names. The second, \texttt{referer\_suspicion},
is zero for a similar reason: every line in both the CSIC and the synthetic
data carries a placeholder ``\,-\,'' Referer, so the feature has no variance to
exploit. In real production traffic, where the Referer header varies from
request to request, the same feature would help detect cross-site request
forgery, hot-linking, and requests arriving from mismatched or suspicious
origins. The third, \texttt{path\_special\_chars}, is near-zero for a different
reason entirely: the information it carries --- the presence of special
characters in the URL path --- is strongly correlated with, and therefore
already captured by, the much stronger entropy, length and encoding features,
so the Random Forest gains almost no additional node purity by splitting on it.
It nevertheless remains a cheap, interpretable corroborator and would carry
more weight on traffic where the attack is delivered through the path segment
rather than the query string.

The same reasoning applies to the small but non-zero features such as
\texttt{signature\_score} ($0.8\%$), \texttt{is\_post} ($0.5\%$),
\texttt{suspicious\_ua\_score} ($0.3\%$) and
\texttt{consecutive\_encoding\_intensity} ($0.2\%$). They contribute little on
this particular corpus, but they encode signals that become important under
different threat mixes: \texttt{suspicious\_ua\_score} would rise against
automated scanners that send anomalous User-Agent strings, and
\texttt{consecutive\_encoding\_intensity} would rise against payloads hidden
behind several layers of obfuscating encoding. Two general points follow.
First, Gini importance is dataset- and correlation-dependent: a feature that is
constant or redundant on this balanced corpus can become decisive on a
different traffic distribution, so retaining it improves the system's
robustness and its ability to generalise to attack types and deployments not
represented here. Second, several of these low-importance features are
consumed by the rule-based Tier~1 and the one-class Tier~3 rather than by the
Random Forest alone, so they still contribute to the overall hybrid even when
the supervised tier does not rely on them. Since all $22$ features are computed
in a single inexpensive pass over each log line, keeping the full set costs
almost nothing while preserving this resilience.

% =====================================================================
\section{Performance of the Hybrid Fusion Strategies}
\label{sec:res-hybrid}

The second experiment evaluates whether fusing the tiers improves on the best
single tier. Table~\ref{tab:hybrid} compares the three tiers against the two
fusion strategies, and Figure~\ref{fig:tierhybrid} plots the same five
configurations against a $90\%$ reference line for each metric.
Figure~\ref{fig:pr} shows the precision--recall trade-off as a scatter plot,
and Figure~\ref{fig:radar} overlays the three strongest configurations on a
radar chart.

\begin{table}[ht]
  \centering
  \caption{Tiers versus hybrid fusion strategies on the evaluation set
           (all values in \%). The best value in each column is in bold.}
  \label{tab:hybrid}
  \renewcommand{\arraystretch}{1.35}
  \setlength{\tabcolsep}{9pt}
  \footnotesize
  \begin{tabular}{|l|r|r|r|r|}
    \hline
    \textbf{Configuration} & \textbf{Prec.} & \textbf{Recall} & \textbf{F1} &
    \textbf{F2} \\
    \hline
    Tier 1: Regex rules & \textbf{99.61} & 24.07 & 38.77 & 28.37 \\
    \hline
    Tier 2: Random Forest & 94.80 & 94.76 & \textbf{94.78} & 94.77 \\
    \hline
    Tier 3: Local Outlier Factor & 88.28 & 95.13 & 91.57 & 93.67 \\
    \hline
    Simple Voting (T1 OR T2 OR T3) & 87.78 & \textbf{97.96} & 92.59 & 95.74 \\
    \hline
    Smart Consensus (proposed) & 91.19 & 97.62 & 94.30 & \textbf{96.26} \\
    \hline
  \end{tabular}
\end{table}

\begin{figure}[ht]
  \centering
  \includegraphics[width=\textwidth]{Figure/model_comparison.png}
  \caption{Tier and hybrid configuration comparison across the four metrics.
           The dashed line marks the $90\%$ reference level. Both fusion
           strategies exceed every single tier on recall and F2.}
  \label{fig:tierhybrid}
\end{figure}

\begin{figure}[ht]
  \centering
  \includegraphics[width=0.75\textwidth]{Figure/precision_vs_recall.png}
  \caption{Precision--recall trade-off. The ideal corner is the top-right.
           The hybrid configurations sit in the high-recall, high-precision
           region, while the Isolation Forest and One-Class SVM occupy the
           high-recall but low-precision region.}
  \label{fig:pr}
\end{figure}

\begin{figure}[ht]
  \centering
  \includegraphics[width=0.7\textwidth]{Figure/radar_comparison.png}
  \caption{Multi-metric radar comparison of the three strongest
           configurations: the Random Forest, the Local Outlier Factor, and
           the smart consensus. The smart consensus dominates on the
           recall and F2 axes.}
  \label{fig:radar}
\end{figure}

The central result of the thesis is visible in Table~\ref{tab:hybrid} and
Figure~\ref{fig:tierhybrid}. The proposed smart consensus attains the highest
F2 score of the entire system at $96.26\%$, exceeding the best single tier
(the Random Forest at $94.77\%$) and every other configuration. It does so by
combining a recall of $97.62\%$ with a precision of $91.19\%$: it catches
about $98$ of every $100$ attacks while keeping roughly nine false alarms per
hundred alerts. Compared with the Random Forest alone, the consensus trades a
small amount of precision (from $94.80\%$ to $91.19\%$) for a substantial gain
in recall (from $94.76\%$ to $97.62\%$); because the F2 metric rewards recall,
this trade is favourable for an intrusion detection system, where missed
attacks are the more dangerous error.

The comparison between the two fusion strategies explains why the consensus
logic, and not mere combination, is responsible for the gain. Simple voting,
which raises an alert whenever any tier fires, reaches the highest recall of
all ($97.96\%$) but its precision drops to $87.78\%$, because it inherits
every false positive of the Local Outlier Factor. The smart consensus, by
requiring that a Tier~3 anomaly be corroborated by at least a mild Random
Forest suspicion before it is trusted, recovers precision to $91.19\%$ while
sacrificing almost no recall (only $0.34$ percentage points). The net effect
is a higher F1 ($94.30\%$ versus $92.59\%$) and a comparable F2, confirming
that the supervised tier's validation of the anomaly tier is the key
mechanism. The precision--recall scatter in Figure~\ref{fig:pr} makes this
trade-off geometric: the hybrid points are pushed towards the desirable
top-right corner, away from the low-precision cluster of the Isolation Forest
and One-Class SVM.

A natural question is how the two probability thresholds that govern the
consensus were chosen: the $0.5$ at which the Random Forest fires on its own,
and the $0.3$ at which it is allowed to confirm a Tier~3 anomaly
(Equation~3.3 of Chapter~3). The two values play different roles and have
different justifications. The first, $0.5$, is not tuned at all: it is the
standard decision boundary of a probabilistic classifier, the point at which
the model estimates an attack to be more likely than a benign request. This is
the Bayes-optimal threshold when the two kinds of misclassification are treated
as equally costly, so the Random Forest raises an alert on its own exactly when
it is more confident of an attack than of a clean request, with no
application-specific calibration. The second value, $0.3$, is a deliberate,
principled relaxation rather than an arbitrary number. It is applied
\emph{only} when the independent anomaly tier has already flagged the request,
and it rests on the observation that two weak but independent signals are
jointly stronger than either alone: once the unsupervised tier has raised a
flag, requiring only a moderate---rather than a majority---supervised suspicion
to confirm it is sufficient. The specific level, $0.3$, embodies the system's
recall-first ($F_2$) priority: it is low enough to recover attacks that the
Random Forest at $0.5$ would have missed but the anomaly tier caught, yet high
enough to still demand non-trivial supervised agreement, which is precisely
what filters out the many false positives the anomaly tier would otherwise
contribute. The thresholds are therefore motivated \emph{theoretically}---$0.5$
from decision theory and $0.3$ from the corroboration principle and the
recall-first objective---and their combined effect is then validated
\emph{empirically} by the comparison in this very section. The two limiting
cases confirm that the value is sensibly placed: had the lower threshold been
raised towards $0.5$, the consensus would collapse back onto the Random Forest
alone and forfeit the recall gain, whereas had it been lowered towards zero, it
would degenerate into the simple-voting baseline and inherit its low precision
of $87.78\%$. The chosen value sits between these extremes, and the ablation
shows it yields the best recall-oriented balance, namely the highest $F_2$
($96.26\%$) at a precision of $91.19\%$.

It is worth noting how these figures relate to the related work surveyed in
Chapter 2. Although a direct numeric comparison is
imperfect because the evaluation protocols and feature sets differ, the
proposed hybrid attains a recall ($97.62\%$) that is higher than the
attack-class recall reported by Pham et al.~\cite{pham2016} for their best
model, while operating directly on raw logs and, unlike the purely supervised
studies, retaining an unsupervised tier capable of flagging novel attacks.
The contribution of this thesis is therefore not a marginal accuracy gain on
a fixed benchmark but a fusion design that is simultaneously
recall-oriented, robust to unseen attacks, and explainable.

% =====================================================================
\section{Error Analysis via Confusion Matrices}
\label{sec:res-confusion}

To understand \emph{where} each configuration makes mistakes, this section
examines the confusion matrices in Figure~\ref{fig:confusion}. Each matrix
reports, for one configuration, the four counts $TN$, $FP$, $FN$, and $TP$;
the off-diagonal cells are the errors, with the bottom-left cell ($FN$) being
the missed attacks and the top-right cell ($FP$) being the false alarms.

\begin{figure}[ht]
  \centering
  \includegraphics[width=\textwidth]{Figure/confusion_matrices.png}
  \caption{Confusion matrices of the eight configurations on the same
           $28{,}767$-line evaluation set. For each panel the bottom-left
           cell counts missed attacks ($FN$) and the top-right cell counts
           false alarms ($FP$); the two diagonal cells are correct decisions.
           The per-panel F1 and F2 values match Tables~\ref{tab:individual}
           and~\ref{tab:hybrid}.}
  \label{fig:confusion}
\end{figure}

The matrices tell a consistent error story, and because they are computed on
the same evaluation set their cell counts can be read directly. The regex-only
panel has an almost empty false-alarm cell (only $13$ false positives) but a
very large missed-attack cell ($10{,}359$ of the $13{,}643$ attacks missed),
the visual signature of high precision and low recall. The Isolation Forest
and One-Class SVM panels show the opposite imbalance: their missed-attack
cells are small (around $900$--$1{,}200$) but their false-alarm cells are
enormous (about $7{,}500$ clean requests wrongly flagged), reflecting their
low precision. The Random Forest panel is well balanced, with small and almost
equal errors on both diagonals ($715$ missed attacks and $709$ false alarms).

The two right-most panels make the effect of the fusion concrete. Simple
voting achieves the fewest missed attacks of all ($278$) but pays with
$1{,}860$ false alarms, because it inherits every Local Outlier Factor false
positive. The smart consensus keeps the missed-attack cell almost as low
($325$ missed attacks, against $665$ for the Local Outlier Factor alone)
while cutting the false alarms to $1{,}286$, well below both simple voting
and the Local Outlier Factor's own $1{,}723$. In other words, requiring the
Random Forest to corroborate a Tier~3 anomaly removes roughly a quarter of the
anomaly tier's false positives at the cost of only a handful of additional
missed attacks. This is exactly the behaviour an intrusion detection system
should exhibit: it errs, when it must, on the side of raising an alert rather
than letting an attack through, and the confusion matrices therefore
corroborate the quantitative conclusion of Section~\ref{sec:res-hybrid} that
the smart consensus offers the best recall-oriented balance.

% =====================================================================
\section{Cross-dataset Evaluation and Dataset Selection}
\label{sec:res-crossdataset}

All results reported so far were obtained on the primary corpus, namely
CSIC~2010 combined with the GitHub attack payloads and synthetic traffic. To
check that the proposed methodology is not tied to a single benchmark, the
entire three-tier pipeline was re-trained and re-evaluated, under an
identical $70/30$ protocol, on two additional public web-attack datasets:
the HttpParams~2015 parameter-value dataset and the ECML/PKDD~2007 challenge
dataset. Table~\ref{tab:crossdataset} reports, for each dataset, the F1 score
of the best supervised model (Random Forest) and the best unsupervised model
(Local Outlier Factor), together with the F1 and F2 of the proposed smart
consensus.

\begin{table}[ht]
  \centering
  \caption{The three-tier system trained and tested on each dataset under the
           same protocol (all values in \%).}
  \label{tab:crossdataset}
  \renewcommand{\arraystretch}{1.35}
  \setlength{\tabcolsep}{8pt}
  \footnotesize
  \begin{tabular}{|l|r|r|r|r|r|}
    \hline
    \textbf{Dataset} & \textbf{Eval size} & \textbf{RF F1} &
    \textbf{LOF F1} & \textbf{Consensus F1} & \textbf{Consensus F2} \\
    \hline
    CSIC 2010 + GitHub (primary) & 28{,}767 & 94.78 & 91.57 & 94.30 & 96.26 \\
    \hline
    HttpParams 2015 & 9{,}321 & 99.73 & 82.25 & 99.53 & 99.61 \\
    \hline
    ECML/PKDD 2007 & 15{,}035 & 83.09 & 57.24 & 73.74 & 79.26 \\
    \hline
  \end{tabular}
\end{table}

The three datasets behave very differently, and this is precisely what
motivates the choice of CSIC~2010 as the primary corpus. On
\textbf{HttpParams~2015} the system reaches almost perfect scores
(Random Forest F1 of $99.73\%$), but this dataset consists of isolated
parameter \emph{values} built from textbook attack payloads, so the decision
boundary between benign and malicious samples is unusually clean; in effect
the dataset makes the system look ``too easy'' and does not reflect the
structure of real HTTP traffic. On \textbf{ECML/PKDD~2007} the scores are the
lowest (Random Forest F1 of $83.09\%$, consensus F2 of $79.26\%$) because its
requests are heavily randomised: paths, parameters, and headers are
synthetic, noisy strings, so even legitimate traffic looks irregular and the
detectors struggle to learn a stable notion of ``normal''. \textbf{CSIC~2010},
by contrast, was generated against a real e-commerce web application, so its
requests have realistic endpoints, parameter names, and value distributions
that closely match the traffic a deployed system would actually see; it
yields strong and, above all, \emph{representative} results
(consensus F2 of $96.26\%$). For these reasons CSIC~2010 combined with the
GitHub payloads was selected as the main dataset of the thesis: it offers the
best balance between high, trustworthy accuracy and fidelity to real-world
web traffic, whereas HttpParams is unrealistically easy and ECML is
unrealistically noisy. Across all three datasets, however, the qualitative
ranking of the tiers is consistent---the Random Forest is the strongest
single model and the smart consensus delivers the highest F2---which confirms
that the proposed methodology generalises and is not an artefact of one
benchmark.

% =====================================================================
\section{Chapter Summary}
\label{sec:ch4summary}

This chapter has evaluated the proposed system in full.
Section~\ref{sec:eval-params} fixed the evaluation parameters: the
precision, recall, F1, and F2 metrics, with F2 as the primary measure, on a
balanced evaluation set of $28{,}767$ lines.
Section~\ref{sec:sim-method} described the simulation method, including the
dataset preparation, the choice of the regex, Logistic Regression, Isolation
Forest, One-Class SVM, and simple-voting baselines, the fixed
hyperparameters, and the deterministic per-sample evaluation procedure. The
results were then reported in three parts. The individual-tier experiment
(Section~\ref{sec:res-individual}) established the Random Forest as the best
supervised model and the Local Outlier Factor as the best unsupervised model,
and confirmed the high-precision, low-recall nature of the regex baseline.
The hybrid experiment (Section~\ref{sec:res-hybrid}) showed that the proposed
smart consensus achieves the system's best F2 score of $96.26\%$, improving on
every single tier and on naive voting by validating anomaly detections with
the supervised model. Finally, the error analysis
(Section~\ref{sec:res-confusion}) confirmed, through the confusion matrices,
that the smart consensus minimises missed attacks while keeping false alarms
under control. Finally, the cross-dataset study
(Section~\ref{sec:res-crossdataset}) re-trained the pipeline on HttpParams
2015 and ECML/PKDD 2007 and justified the choice of CSIC~2010 as the primary
corpus, showing that the tier ranking is consistent across datasets while
CSIC offers the most realistic and trustworthy results. Together, these
results demonstrate that the three-tier hybrid design meets its objective of
recall-oriented, robust, and explainable web intrusion detection. The next
chapter draws conclusions and outlines directions for future work.


\end{document}